% !TEX program = pdflatex
% ---------------------------------------------------------------------------
% Global Shocks, Contaminated Controls (M1 rev.2)
% English manuscript -- reframed as a failure-diagnosis paper
% Compile: latexmk -pdf paper.tex   (bibtex handled automatically)
% ---------------------------------------------------------------------------
\documentclass[11pt,a4paper]{article}

\usepackage[margin=1in]{geometry}
\usepackage{amsmath,amssymb}
\usepackage{booktabs}
\usepackage{graphicx}
\usepackage{multirow}
\usepackage[round]{natbib}
\PassOptionsToPackage{hyphens}{url}% allow long URLs in bibliography to break
\usepackage[colorlinks=true,linkcolor=blue,citecolor=blue,urlcolor=blue]{hyperref}

\renewcommand{\arraystretch}{1.15}

% Author info -- adjust before submission
\newcommand{\authorname}{Changhao Li}
\newcommand{\authoremail}{(email to be inserted on submission)}
\newcommand{\authoraffil}{The Hong Kong Polytechnic University}

\title{\textbf{Global Shocks, Contaminated Controls:\\
Diagnosing Counterfactual Failure in Trump's 2018 and 2025 Tariff Events\thanks{%
Working paper, September 2026 (rev.~2). Replication code and data acquisition
scripts are available at the companion repository
\texttt{quant\_strategy/notebooks/paper\_tariff/}; all data are from Yahoo
Finance via \texttt{yfinance}. Comments welcome.}}}

\author{\authorname\thanks{\authoremail. \ \authoraffil}}
\date{September 2026}

\begin{document}
\maketitle

\begin{abstract}
\noindent
We study what counterfactual methods can and cannot deliver when a policy
shock treats every candidate control at once. On Donald Trump's April 2025
``Liberation Day'' the S\&P~500 fell 12.1\% over the four trading sessions
from April 3 to April 8 before rebounding on the April~9 pause; every
candidate control---U.S.\ sector ETFs
(slices of the treated index) and overseas indices (all sharply down over
the same window)---was itself treated. We diagnose the consequences rather
than report effects. Placebo tests have no power by construction when donors
are homogeneous slices of the treated unit; synthetic weights concentrate on
the single best pre-period tracker; and counterfactual gaps order
monotonically with how tightly the counterfactual is bound to its
donors---near zero under the tightest binding, approaching the raw return as
binding loosens. The spread across specifications therefore measures control
construction, not treatment effects, and the various ``gaps'' are different
estimands rather than one quantity with uncertainty. A mediation-style
attribution of the crash to a VIX ``fear channel'' fails on two counts:
sequential ignorability is violated, and VIX is the implied volatility of
S\&P~500 options, so the exercise measures mechanical contemporaneous
covariation, not a channel. What survives without any donor is descriptive:
VIX rose about 140\% from its pre-announcement close to its April 8 peak
in the 2025 crash and decayed
with a one-day half-life, while the comparable 2018 node produced no
immediate VIX reaction---a contrast consistent with, but not identified
as, market learning. We distill these diagnostics into a checklist for
synthetic-control applications to global financial shocks.

\medskip
\noindent\textbf{Keywords:} synthetic control; contaminated controls; event
study; tariff shock; identification failure; volatility \\
\noindent\textbf{JEL:} G14, G18, F13, C54

\end{abstract}

\newpage
\tableofcontents
\newpage

% ===========================================================================
\section{Introduction}\label{sec:intro}

The tariff policies of Trump's two presidencies constitute the largest
American trade intervention in a century, and they arrived in two delivery
mechanisms: the gradual 2018--19 escalation through Section~232 and four
Section~301 lists (roughly \$370 billion of goods cumulatively), and the
shock of April 2, 2025 (``Liberation Day''), when a universal baseline plus
reciprocal rates were announced against virtually all trading partners at
once. A natural question for market participants and policy evaluators
alike: what did these shocks do to the S\&P~500 relative to the
counterfactual of no tariff war?

This paper's answer is that the question cannot be answered by the standard
toolkit in this setting, and that the reasons are instructive. The synthetic
control method of \citet{abadie2010} and its Bayesian-regression cousin
\citep{brodersen2015} require controls untouched by the intervention. A
global tariff shock treats every candidate control simultaneously: U.S.\
sector ETFs are components of the treated index; every overseas developed
market fell sharply over the same event window. Once this is recognized,
the interesting research question shifts from ``what is the effect?'' to
``what do these estimators produce when their identifying assumption fails
in a known direction, and how can a researcher tell?''

We therefore reframe the exercise as a \emph{failure diagnosis with a
surviving descriptive core}. Across two tariff rounds, two donor families,
and a battery of estimators (convex SCM, augmented SCM with ridge
regularization, BSTS regression, linear and neural-network maps), we
document four systematic facts and one set of donor-free findings:

\begin{enumerate}
  \item \textbf{No power by construction.} When donors are homogeneous
  slices of the treated unit, the synthetic control reconstructs the
  treated series and the gap tends to zero regardless of the intervention;
  permutation placebos are then uninformative ($p=100\%$: the realized gap
  is the \emph{smallest} in absolute value among all units).
  \item \textbf{Weight concentration.} In the overseas pool, SCM loads
  0.76 on a single index (Canada) and parallel trends fail outright;
  enlarging the pool does not dilute the concentration.
  \item \textbf{Gaps order with donor-binding tightness.} Across an
  augmented-SCM ridge sweep, the gap moves monotonically from the convex
  bound (near zero) to the raw return itself as regularization loosens the
  donors' grip on the counterfactual. A statistic that orders with a
  tuning parameter of the control construction is not measuring the
  treatment.
  \item \textbf{Different estimands masquerading as one interval.} The
  spread of long-window statistics across specifications (roughly two
  orders of magnitude) is sometimes presented as an ``uncertainty
  boundary.'' It is better described as a collection of different
  estimands---S\&P~500 minus sector-synthetic, S\&P~500 minus an
  overseas-predicted path, nonlinear prediction errors---none of which is a
  tariff effect.
\end{enumerate}

What survives without any control group is descriptive: the volatility
signature. VIX rose from its eve-of-announcement close of 21.8 (April 1)
to 52.3 at the April 8 trough of the 2025
crash (+140\%) and decayed with a one-day half-life after the April~9
pause; the comparable 2018 node (List~1) produced no immediate VIX
reaction at all. We present this contrast as a descriptive fact and
explicitly decline the causal reading (``markets learned'') because
treatment magnitude, surprise, and macro background all differ across
rounds.

A further negative result concerns a tempting exercise: decomposing the
crash into a VIX-mediated ``fear channel'' and a residual direct effect, in
the style of \citet{imai2010}. The exercise fails identification twice.
Sequential ignorability is violated (VIX responds to the equity decline
itself), and---more fundamentally---VIX is the implied volatility of
S\&P~500 options: regressing index returns on $\Delta$VIX is a mechanical
contemporaneous covariation (the leverage/volatility-feedback relation),
not the measurement of a channel. We report the decomposition only as a
descriptive covariation account, with its fragilities stated.

\paragraph{Contributions.}
\begin{enumerate}
  \item \textbf{A diagnosis protocol} for counterfactual methods under
  globally contaminated controls: constructive no-power placebos,
  weight-concentration checks, and a binding-tightness ordering
  (augmented-SCM ridge sweep) that reveals whether a gap tracks the
  control construction rather than the treatment---demonstrated on the
  two-round tariff design and calibrated on a clean-control positive
  control (California Proposition~99).
  \item \textbf{Donor-free descriptive volatility facts} for the two
  rounds, with three competing interpretations (learning; shock magnitude
  and surprise; background conditions) stated and left unresolved.
  \item \textbf{Explicitly non-causal relative accounting}: S\&P~500
  relative to an overseas-predicted path (BSTS, with its
  contaminated-covariate assumption disclosed), and a beta decomposition of
  the sector cross-section showing that raw ``which sector fell most''
  rankings are dominated by market beta. These are descriptions of relative
  performance, not tariff effects.
\end{enumerate}

\paragraph{What we do not claim.} We do not estimate the effect of tariffs
on the S\&P~500; we do not interpret any SCM/BSTS gap as causal; we do not
claim a fear ``channel'' share; and we do not claim evidence of market
learning. The paper's value is in showing precisely why each of these
claims would fail, and what can be said instead.

The remainder proceeds as follows. Section~\ref{sec:background} lays out
the institutional timeline, including the 2026 litigation aftermath.
Section~\ref{sec:data} describes the data. Section~\ref{sec:method} sets
out the estimators and the diagnostic uses to which we put them.
Section~\ref{sec:results} presents the diagnosis. Section~\ref{sec:protocol}
distills the diagnostics into a pre-reporting protocol and demonstrates it
on a clean-control positive control. Section~\ref{sec:discussion} discusses
interpretation and limitations, and Section~\ref{sec:conclusion} concludes.

% ===========================================================================
\section{Institutional Background}\label{sec:background}

\subsection{Two tariff rounds}

Table~\ref{tab:timeline} summarizes the key dates. The essential difference
is the delivery mechanism: \textbf{2018 was gradual}---tariffs arrived in
batches over eighteen months and markets priced each step as it came;
\textbf{2025 was a shock}---a single announcement covering all major
partners at once, followed by a crash and an equally fast partial reversal.
Our focal events are 2018-07-06 (List~1, the largest single node of round
one) and 2025-04-02 (Liberation Day).

\begin{table}[htbp]
\centering
\caption{Tariff timeline, 2018--2026 (selected events; compiled from
\citealp{bownkolb2025} and Federal Register and court records)}
\label{tab:timeline}
\small
\begin{tabular}{@{}llp{9.0cm}@{}}
\toprule
Round & Date & Event \\
\midrule
2018 & Mar 08 & Section 232 steel (25\%) and aluminum (10\%) tariffs signed \\
2018 & Mar 22 & Section 301 China memorandum signed \\
2018 & Jul 06 & \textbf{List 1 effective, 25\% (\$34B)} --- round-1 focal event \\
2018 & Aug 23 & List 2 effective, 25\% (\$16B) \\
2018 & Dec 01 & G20 Buenos Aires truce (90 days) \\
2019 & May 10 & List 3 raised 10\%$\to$25\% (\$200B) \\
2019 & Sep 01 & List 4A effective, 15\% (\$112B) \\
\midrule
2025 & Feb 01 & China +10\%, Canada/Mexico +25\% (IEEPA/fentanyl) \\
2025 & Apr 02 & \textbf{Liberation Day reciprocal tariffs} --- round-2 focal event \\
2025 & Apr 09 & 90-day pause on reciprocals; China rate raised to 125\% \\
\midrule
2026 & Feb 20 & \emph{Learning Resources v.\ Trump}: Supreme Court (6--3,
Roberts, C.J.) holds IEEPA does not authorize the tariffs \\
2026 & Feb 24 & Government pivots to Section 122 (Trade Act of 1974): 10\%
global surcharge (statutory cap 15\%, sunset 150 days) \\
2026 & May 07 & Court of International Trade (2--1) holds Section 122
unlawful; CAFC administrative stay; CBP keeps collecting \\
2026 & Jul 24 & Section 122 statutory expiry; 60 Section 301 investigations
(initiated Mar 12) take over as successor authority \\
\bottomrule
\end{tabular}
\end{table}

\subsection{The 2026 legal aftermath}

The legal foundation of the 2025 tariffs collapsed within a year of the
event we study: the Supreme Court held in \emph{Learning Resources, Inc.\
v.\ Trump} that IEEPA does not authorize presidential tariff imposition
\citep{learningresources2026}; the administration immediately pivoted to a
Section~122 surcharge \citep{section122tariffs2026}, which the Court of
International Trade invalidated \citep{citsection1222026} before its
150-day statutory expiry, with Section~301 investigations
\citep{section301investigations2026} standing by as the third legal basis
in twelve months. This relay matters for our interpretation: it is direct
evidence that tariff policy in this period was a moving target whose legal
durability was unknown at each moment---economic policy uncertainty in the
classic sense of \citet{baker2016}---and hence background condition, not
object, of our analysis.

% ===========================================================================
\section{Data}\label{sec:data}

All price data are daily closes from Yahoo Finance via \texttt{yfinance},
2017-01-01 through 2025-12-30. Table~\ref{tab:data} lists the panels.

\begin{table}[htbp]
\centering
\caption{Data panels}
\label{tab:data}
\small
\begin{tabular}{@{}lp{5.4cm}p{7.4cm}@{}}
\toprule
Panel & Instruments & Role \\
\midrule
us\_index & \^{}GSPC \^{}DJI \^{}IXIC & Treated unit (S\&P~500) \\
sector\_etf & XLY \ldots{} XLRE (11 ETFs) & Candidate donor pool 1 \\
overseas\_index & 15 developed markets\textsuperscript{a} & Candidate donor
pool 2 (extended 2026-07-12) \\
macro & \^{}VIX, GC=F, \^{}TNX, DX-Y.NYB & Covariates \\
\bottomrule
\end{tabular}

\vspace{2pt}
\noindent\textsuperscript{a}\small \^{}N225, \^{}GDAXI, \^{}FTSE, \^{}FCHI,
\^{}AXJO, \^{}GSPTSE, \^{}STOXX50E, \^{}SSMI, \^{}N100, VGK, EZU, EWJ, EPP,
EWA, EWC; EEM (emerging) and HSI (China-linked) excluded.
\normalsize
\end{table}

\paragraph{Limitations of the data.} \texttt{yfinance} covers currently
listed tickers; sector-ETF survivorship is mild but nonzero. XLC (listed
2018-06) and XLRE (2018-10) lack pre-period data for the 2018 round, so the
two-round comparison (Section~\ref{sec:tworound}) uses the 9 common donors.
Sector China-exposure figures (Section~\ref{sec:exposure}) are PIIE/BEA-based
proxies. VIX futures term-structure data are not reliably available from
free sources, so all volatility statements use spot VIX. A GDELT news panel
(1{,}416 items over 2025-04-01--15) was collected for an auxiliary sentiment
dimension whose finance-specific tone proved weak (contemporaneous
correlation with S\&P~500 returns $+0.06$); it is excluded from the
analysis and documented in Appendix~\ref{app:repro} only.

% ===========================================================================
\section{Methods and Their Diagnostic Uses}\label{sec:method}

The estimators below are standard; what is unusual is the use we put them
to. Because the central assumption of every one of them (untreated
controls) fails in a known direction, we treat each estimator's output as a
\emph{diagnostic statistic} whose behavior under contamination we document,
rather than as an effect estimate.

\subsection{Synthetic control}\label{sec:scm}

Following \citet{abadie2010}, with pre-treatment period $T_0$ we solve for
donor weights $w$:
\begin{equation}
\min_{w} \;\sum_{v} V_v\Bigl(X_{1v}-\sum_{j} w_j X_{jv}\Bigr)^{2}
+\bigl\lVert Z_1 - Z_0 w\bigr\rVert^{2}
\qquad \text{s.t.}\;\; w_j\ge 0,\;\; \textstyle\sum_j w_j = 1 ,
\label{eq:scm}
\end{equation}
where $V_v$ weights the predictors. We report both equal $V$ and the
data-driven choice of \citet{abadie2021} (\texttt{pysyncon}), because
Section~\ref{sec:vsens} shows the gap is sensitive to it. Weights are
estimated with \texttt{cvxpy} (CLARABEL) and cross-checked against
\texttt{pysyncon}. Gaps are computed on cumulative log returns.

\subsection{Inference tools read as power diagnostics}\label{sec:inference}

The \emph{permutation placebo} \citep{abadie2021} iteratively treats each
donor as the treated unit; the permutation $p$-value is the share of
placebo gaps at least as extreme as the realized one. Under donor
homogeneity this test has no power by construction, and a $p$-value of
100\% is itself the diagnostic finding. \emph{Leave-one-out} (LOO)
re-estimation exposes dependence on any single donor. \emph{In-time
placebos} move the intervention date into the pre-period and locate the
realized statistic in the resulting null distribution; we report these
quantiles as calibration statistics, not as significance tests of an
effect.

\subsection{Bayesian structural time series as relative accounting}
\label{sec:bsts}

\citet{brodersen2015} fit a state-space model (local level plus covariate
regression) on the pre-period and forecast a counterfactual. We implement
the frequency-domain equivalent with \texttt{statsmodels}
\texttt{UnobservedComponents}. Two assumptions fail here: the covariates
(overseas indices) are treated, and---unlike convex SCM---the regression
extrapolates without simplex constraints, so contamination passes through
coefficients rather than being averaged. We therefore read its output
exclusively as \emph{relative-performance accounting}: where the S\&P~500
stood relative to an overseas-predicted path over the same window.

\subsection{Event study and its limits}\label{sec:es-method}

Following \citet{mackinlay1997}, the market model
$R_{it} = \alpha_i + \beta_i R_{mt} + \varepsilon_{it}$ is estimated over
an estimation window and abnormal returns cumulated over the event window.
In a market-wide shock, however, the market factor is itself the treated
variable and contains each sector's own return; residuals therefore
describe \emph{relative} performance, not causal abnormal returns. We use
the event study for one descriptive purpose: decomposing the sector
cross-section into a beta component and a residual component
(Section~\ref{sec:cross-section}).

\subsection{Volatility-covariation accounting (not mediation)}
\label{sec:covar-method}

For any return series we estimate, over the pre-period,
\begin{equation}
R_{t} = \alpha + \beta\, \Delta VIX_{t} + \varepsilon_{t},
\qquad \Delta VIX_t = \text{VIX log return},
\label{eq:vixreg}
\end{equation}
and use it two ways. First, an \emph{orthogonalized SCM}: treated and every
donor are purged with their own pre-estimated $\beta$ (out of sample over
the event window), and the SCM re-run on the purified paths. Second, a
\emph{descriptive decomposition} of the event-window cumulative return:
$\sum_{t\in E} R_t = \hat\beta \sum_{t\in E}\Delta VIX_t
+ \sum_{t\in E}\hat\varepsilon_t$, the first term measuring contemporaneous
covariation with the volatility state.

Three properties of this exercise must be stated before any number is
read. (i) $\Delta$VIX is contemporaneous, and VIX is the implied volatility
of S\&P~500 options: the return--$\Delta$VIX relation embeds the
leverage/volatility-feedback mechanics, so covariation shares are close to
re-descriptions of the decline itself. (ii) $\hat\beta$ is estimated on a
calm pre-period and extrapolated across a regime shift
(Section~\ref{sec:garch}). (iii) Sequential ignorability, required for
causal mediation \citep{imai2010}, fails outright. We therefore do not use
the language of channels, mediation, or causal shares anywhere in the
results.

\subsection{GARCH and VIX paths (donor-free)}

For each round we fit a GARCH(1,1) with an event dummy in the variance
equation,
\begin{equation}
\sigma_t^2 = \omega + \alpha\,\varepsilon_{t-1}^2 + \beta\,\sigma_{t-1}^2
+ \gamma\, D_t ,
\label{eq:garch}
\end{equation}
and, more importantly, read the shape of the VIX path (level change, time
to peak, half-life). Neither statistic uses a control group; both are
descriptive.

\subsection{Augmented SCM as a binding-tightness dial}\label{sec:ascm-method}

The convex hull constraint of equation~\eqref{eq:scm} confines the
counterfactual to the donor simplex. The augmented SCM of
\citet{benmichael2021} (see also \citealp{doudchenko2016}) replaces convex
weights with ridge regression, anchored at the convex-SCM benchmark and
level-anchored at the event date, so results remain comparable to SCM cumlog
gaps. We use the ridge penalty sweep $\alpha\in[0.01,100]$ as a dial that
moves the counterfactual continuously from ``fully bound to donors''
(convex) to ``essentially unbound'' (level line). If the gap orders
monotonically along this dial, it tracks the control construction; the
script \texttt{ascm\_sweep.py} reproduces Table~\ref{tab:ascm}.

\subsection{Neural-network maps as extrapolation-error probes}

MLP and LSTM models map donor returns plus VIX to S\&P~500 returns
(pre-period training, event-period out-of-sample application; 5 seeds;
window lengths 3/5/10). Their outputs are \emph{prediction errors of a
nonlinear pre-period map}---structurally contaminated in the same way as
SCM---and we report them only to document the magnitude and instability of
nonlinear extrapolation error, never as gaps.

% ===========================================================================
\section{Results: The Diagnosis}\label{sec:results}

\subsection{The Liberation Day window}\label{sec:narrative}

The reciprocal tariffs were announced after the close on 2025-04-02
(close: 5{,}671). Over the following four trading sessions (April 3, 4, 7,
and 8) the S\&P~500 fell to a trough close of 4{,}983 on April 8
($-12.1\%$); the April~9 pause triggered a
$+9.52\%$ single-day rebound. The sector cross-section ranged from XLE
(energy) $-18.8\%$ to XLP (staples) $-6.4\%$. Every overseas index in the
donor pool fell over the same window ($-7.6\%$ to $-13.3\%$). Two
properties of this window drive everything that follows: the shock was
global (no market was untreated), and the domestic ``controls'' are
components of the treated index.

\subsection{Sector SCM: homogeneity by construction}\label{sec:vsens}

Over the pre-period 2024-06--2025-04, the equal-$V$ SCM loads XLK 0.36,
XLY 0.18, XLV 0.17, XLF 0.09 (top four: 79.8\%); the pre-trend statistic
is $p=0.957$ (Newey--West)---though, as Section~\ref{sec:placebo} shows,
this design has no power, so a high pre-trend $p$-value is uninformative
rather than reassuring. Realized gaps: short window (4/2$\to$4/8)
$-0.04\%$; long window (4/2$\to$6/30) $+0.68\%$. The near-zero gaps are
not evidence that tariffs were harmless; they are what this construction
must produce when donors are slices of the treated index---the synthetic
control approximately \emph{reconstructs} the S\&P~500, weight-fitting
details aside. Switching to data-driven $V$
\citep{abadie2021} disperses the weights and moves the long-window gap to
$+7.45\%$: an 11$\times$ spread from the equal-$V$ benchmark attributable
to the weight matrix alone. A statistic that moves by an order of
magnitude with a defensible nuisance choice is not measuring the treatment.

\subsection{Placebo and leave-one-out: no power by construction}
\label{sec:placebo}

The permutation test returns $p=100\%$: the realized short-window gap
($-0.04\%$) has the smallest absolute value among all 12 units, against
placebo gaps ranging $-6.24\%$ to $+4.69\%$. This is the constructive
no-power result announced in the introduction: with homogeneous donors,
synthetic $\approx$ treated and the gap is pinned near zero
\emph{irrespective of the intervention}; the test cannot reject anything,
ever, in this design. One caveat on interpretation: a smallest-of-all
realized gap is observationally consistent with two readings---the
homogeneity artifact we emphasize, or genuine treatment heterogeneity in
which the cap-weighted, globally diversified index is simply the least
tariff-exposed unit in the pool while concentrated sector funds absorb
larger shocks. Our design cannot separate the two, and nothing downstream
depends on choosing between them. LOO gaps span $[-0.49\%, +0.69\%]$%
---stable, for the same reason: every donor subset still reconstructs the
index. Stability here is a symptom of the design, not a robustness
property of an estimate.

\subsection{Overseas controls: SCM unusable, BSTS as relative accounting}
\label{sec:overseas}

With the overseas pool at 15 developed markets, SCM weights concentrate on
Canada (0.76) and Australia (0.21) --- tickers \^{}GSPTSE and \^{}AXJO ---
and pre-trends fail outright ($p=0.000$); the long-window statistic
($+2.43\%$) is uninterpretable and we discard it. Extending the pool did not dilute
concentration (an earlier pool put 0.99 on EEM): in any overseas family
there is always one index that tracked the S\&P~500 best pre-event, and
convex SCM gives it the weight.

The BSTS regression of the S\&P~500 on the equal-weighted overseas basket
is retained only as relative accounting: the S\&P~500 closed the short
window $3.54\%$ below its overseas-predicted path and the long window
$4.35\%$ above it (posterior tail-area statistics 60\% and 8.2\%; in-time
placebo quantiles 10\% and 60\%). These numbers describe where the U.S.\
market stood relative to a (treated) world basket---a descriptive
comparison across contaminated markets, in which the U.S.\ policy reversal
of April~9 plausibly plays a role. They are not tariff effects, and we do
not pool them with SCM statistics (Section~\ref{sec:map}).

\subsection{Two rounds: realized reactions versus constructed gaps}
\label{sec:tworound}

Table~\ref{tab:tworound} compares the rounds on the 9 common donors. The
windows deliberately follow each round's natural rhythm---2018: short =
event to $+2$ months (2018-09-06), long = to year-end; 2025: short =
event to $+4$ sessions (2025-04-08), long = to June 30---because the
delivery mechanisms differ; the notebook records this as design, not
oversight. Because the realized cells therefore confound shock with
window length, we re-ran both rounds at \emph{aligned} horizons of $+4$,
$+20$, and $+60$ sessions (\texttt{aligned\_windows.py}). Three readings
follow. (i)~At $+4$ sessions the crash contrast is intact: $+1.4\%$
(2018) versus $-12.1\%$ (2025). (ii)~By $+20$ sessions it has largely
dissipated ($+2.9\%$ versus $-1.2\%$), and at $+60$ both rounds are
positive ($+6.0\%$ versus $+9.4\%$): the dramatic long-window contrast in
Table~\ref{tab:tworound} ($-9.2\%$ versus $+9.4\%$) is
manufactured by unequal windows---the 2018 long cell runs six months and
absorbs the Q4 2018 drawdown, while at the aligned three-month point
(October 1) the 2018 return stood at $+6.0\%$. (iii)~In every aligned
cell the SCM statistic stays below 1\% (range $-0.09\%$ to $+0.91\%$),
consistent with the construction story of Section~\ref{sec:vsens}. The
price cells---realized or synthetic---are not the carrier of the
between-round difference; anything substantive lives in the volatility
signature (Section~\ref{sec:garch}). Note also that the 2025 short-window
statistic here ($-0.09\%$) uses the 9-donor pool and therefore differs
slightly from the 11-donor figure of Section~\ref{sec:vsens}
($-0.04\%$); both are reported, neither is an effect.

\begin{table}[htbp]
\centering
\caption{Two-round comparison (9 common donors; windows follow each
round's rhythm, see text; descriptive---no cell is an effect)}
\label{tab:tworound}
\small
\setlength{\tabcolsep}{4pt}
\begin{tabular}{@{}lcccc@{}}
\toprule
Round & Realized (short / long) & SCM statistic (short / long) & MSPE &
Pre-trend \\
\midrule
2018--19 (List~1) & $+4.3\%$ / $-9.2\%$ & $+0.03\%$ / $-0.69\%$ & 7.2
& $p=0.727$ \\
2025 (Liberation Day) & $-12.1\%$ / $+9.4\%$ & $-0.09\%$ / $+0.91\%$ & 4.7
& $p=0.775$ \\
\bottomrule
\end{tabular}
\end{table}

\subsection{Volatility signatures: the donor-free facts}\label{sec:garch}

These statistics use no control group and are therefore the only findings
in this paper that do not inherit contamination; they are, accordingly,
descriptive statements about realized paths (Figure~\ref{fig:vixpaths}):

\begin{figure}[htbp]
\centering
\includegraphics[width=0.92\textwidth]{fig_vix_paths.png}
\caption{VIX around each focal node, $\pm$60 trading sessions, calendar
axes (a shared relative-event axis would visually compress 2018's null
reaction against 2025's spike). Upper: List~1 (2018-07-06)---the eve
close (15.0) is never exceeded afterward (window peak 14.9). Lower:
Liberation Day (2025-04-02)---eve close 21.8, peak 52.3 on April 8
(+140\%), half-life about one day after the April~9 pause.
Generated by \texttt{make\_vix\_figure.py}.}
\label{fig:vixpaths}
\end{figure}

\begin{itemize}
  \item \textbf{2025 Liberation Day:} VIX closed at 21.8 on the eve of
  the announcement (April 1) and peaked at 52.3 on April 8, four sessions
  into the crash ($+140\%$), then a half-life of about one day---the
  April~9 pause unwound the spike almost immediately.
  \item \textbf{2018 List 1 (Jul 06):} VIX closed at 15.0 on the eve of
  the node (July 5) and never exceeded that level over the following 60
  sessions (window peak 14.9)---no immediate reaction at the node; VIX
  climbed only as nodes
  accumulated into Q4.
\end{itemize}

The GARCH(1,1) fits are secondary: both rounds show high persistence
($\alpha+\beta>0.96$) over $\pm$120-day windows, and the event-dummy
regression is coarse ($\gamma_{18}$ significantly negative, $\gamma_{25}$
insignificant---an ordering built on one insignificant coefficient is not
worth interpreting). We rest nothing on the parameters; the path shapes
carry the content. Three interpretations of the contrast
are live---market learning across rounds; larger and more surprising
treatment in 2025; different rate/inflation/AI-cycle backgrounds---and our
design cannot separate them (Section~\ref{sec:learning}).

\subsection{Volatility-covariation accounting}\label{sec:covar}

Applying equation~\eqref{eq:vixreg}: the treated pre-period coefficient is
$\hat\beta=-0.079$. The orthogonalized SCM (treated and all donors purged,
both with pre-estimated coefficients) moves the statistics only slightly:
short $-0.04\%\to-0.09\%$, long $+0.68\%\to+0.86\%$
(Table~\ref{tab:covar})---the contemporaneous covariation largely cancels
in the treated-minus-synthetic difference because it hits both sides. The
descriptive decomposition of the event-window return attributes a
$-6.85$-point fitted covariation component to the short-window cumulative
move of $-12.27\%$ (residual component $-5.41$; a covariation share of
roughly 56\%; long window $+0.79$ of $+7.64$, residual $+6.85$). Given the three properties stated in
Section~\ref{sec:covar-method}---mechanical contemporaneous link, regime
extrapolation of $\hat\beta$, failed ignorability---this share should be
read as a descriptive statistic about how tightly the decline co-moved
with the options market's repricing of volatility, nothing more. It is not
a channel, not causal, and its precision is illusory; we flag that any
study regressing index returns on their own implied volatility will
``recover'' a large share by construction.

\begin{table}[htbp]
\centering
\caption{SCM statistic before and after VIX-orthogonalization (both
treated and donors purged with pre-period coefficients)}
\label{tab:covar}
\small
\begin{tabular}{@{}lccc@{}}
\toprule
Window & Raw SCM statistic & Purified & Difference \\
\midrule
Short (4/2$\to$4/8) & $-0.04\%$ & $-0.09\%$ & $-0.05$pp \\
Long (4/2$\to$6/30) & $+0.68\%$ & $+0.86\%$ & $+0.18$pp \\
\bottomrule
\end{tabular}
\end{table}

\subsection{Sector cross-section: beta dominates raw rankings}
\label{sec:cross-section}

Over 4/2$\to$4/8 the raw sector ranking (XLE $-18.8\%$ worst) and the
market-model residual ranking differ sharply: on residuals, XLE stands at
$-15.16$ points and XLK at $+5.19$ points, while XLK's raw return was the
worst ($-14.22\%$) against a fitted $\beta$ of 1.59. The cross-sectional
correlation between raw returns and residual ranks is only 0.24: raw
``which sector fell most'' league tables are dominated by market beta
rather than sector-specific exposure. We stop the claim there. Because the
market factor is the treated index and contains the sectors themselves,
residuals cannot be read as causal ``abnormal'' performance; the
$+5.19$-point residual for XLK reflects both its genuine relative strength
during the rebound and the mechanical absorption of its own weight in the
market factor. The usable output is the descriptive decomposition, not a
label of winners and losers.

\subsection{The XLK mechanism}\label{sec:exxlk}

If the near-zero sector-SCM statistic is a homogeneity artifact, removing
the dominant homogeneous donor should break it. It does: excluding XLK
moves the long-window statistic from $+0.68\%$ to $+7.51\%$ (XLK's
contribution $-6.83$ points; post-exclusion weights XLY 0.48, XLV 0.31,
then XLB/XLE; pre-trend $p=0.207$). The mechanism is transparent: with
XLK in the pool the synthetic control inherits the post-event tech-led
rebound and tracks the treated unit; without XLK it misses that rally.
The data-driven-$V$ result of Section~\ref{sec:vsens} ($+7.45\%$) works
through the same channel---both stop matching the pre-period AI rally
with XLK---so the two numbers are one mechanism observed twice, not
independent corroboration.

\subsection{Augmented SCM: gaps order with binding tightness}
\label{sec:ascm}

Table~\ref{tab:ascm} reports the ridge sweep. As $\alpha\to0$ the ridge
approaches unconstrained donor fitting and the statistic approaches the
convex SCM; as $\alpha\to\infty$ coefficients shrink to zero, the
``counterfactual'' degenerates to an anchored level line, and the statistic
converges to the raw return itself ($\alpha=100$: $-11.88\%$ vs.\ realized
$-12.1\%$ short; $+9.40\%$ vs.\ $+9.4\%$ long). At $\alpha=1$ the ridge
coefficients sum to 0.574 (vs.\ 1.0 under convexity): the construction is
part donor, part extrapolation.

\begin{table}[htbp]
\centering
\caption{Augmented SCM: statistic by ridge penalty $\alpha$ (cumlog
metric, event-date anchored; \texttt{ascm\_sweep.py})}
\label{tab:ascm}
\small
\begin{tabular}{@{}lcccccc@{}}
\toprule
Window & SCM (convex) & $\alpha{=}0.01$ & $\alpha{=}0.1$ & $\alpha{=}1$ &
$\alpha{=}10$ & $\alpha{=}100$ \\
\midrule
Short (4/2$\to$4/8) & $-0.04\%$ & $-1.13\%$ & $-2.09\%$ & $-5.06\%$ &
$-9.56\%$ & $-11.88\%$ \\
Long (4/2$\to$6/30) & $+0.68\%$ & $+0.32\%$ & $+1.37\%$ & $+4.19\%$ &
$+7.75\%$ & $+9.40\%$ \\
\bottomrule
\end{tabular}
\end{table}

The monotone ordering is the finding. The statistic sweeps continuously
from the convex bound to the raw return as the dial loosens the donors'
grip: it is a function of the control construction, and there is no
setting of $\alpha$ at which it can be claimed as a treatment effect. The
near-zero convex value is not ``tariffs were harmless''; it is ``under the
tightest binding the design cannot see the shock.''

\subsection{Nonlinear extrapolation error}\label{sec:dl}

The neural maps produce large, unstable numbers: MLP $+10.29\% \pm 7.02$pp
and LSTM $+16.18\% \pm 9.39$pp across 5 seeds (a linear map
gives $+0.16\%$); window sensitivity runs $+15\%$, $+18\%$, and $+21\%$
for $w=3$, $5$, $10$, and single seeds reach $+35\%$. Validation/train
loss ratios
near 1.2 over a more volatile validation segment confirm overfitting. We
draw no inference from these magnitudes beyond what they demonstrate: a
nonlinear pre-period map extrapolated into a regime shift yields errors
two orders of magnitude above the linear map's, unstable across seeds.
These are prediction errors, not counterfactual gaps, and they are
excluded from Table~\ref{tab:map}.

\subsection{Specification map}\label{sec:map}

Table~\ref{tab:map} collects the long-window statistics grouped by what
each actually measures. The grouping is the point: rows within the first
block are one nominal estimand fractured by construction choices; the
second block is a different estimand (relative to a treated world
basket); the third block is not a gap at all. Presenting such a
collection as an ``uncertainty interval'' over one quantity would be a
category error.

\begin{table}[htbp]
\centering
\caption{Specification map: long-window statistics by estimand. No row is
a tariff effect; within-block spread orders with donor-binding tightness
(Sections~\ref{sec:vsens}--\ref{sec:ascm}).}
\label{tab:map}
\small
\begin{tabular}{@{}llc@{}}
\toprule
Estimand & Specification & Long-window statistic \\
\midrule
\multirow{6}{*}{\begin{tabular}[c]{@{}l@{}}S\&P minus\\sector-synthetic\end{tabular}}
 & Convex SCM, equal $V$ & $+0.68\%$ \\
 & Convex SCM, data-driven $V$ & $+7.45\%$ \\
 & Exclude XLK & $+7.51\%$ \\
 & VIX-purified (equal $V$) & $+0.86\%$ \\
 & Linear donor map & $+0.16\%$ \\
 & Augmented SCM, $\alpha\in[0.01,100]$ & $+0.32\%$ to $+9.40\%$ \\
\midrule
\multirow{2}{*}{\begin{tabular}[c]{@{}l@{}}S\&P minus\\overseas-predicted path\end{tabular}}
 & Overseas SCM & $+2.43\%$ (discarded, pre-trend failure) \\
 & BSTS regression & $+4.35\%$ (relative accounting) \\
\midrule
\multirow{2}{*}{\begin{tabular}[c]{@{}l@{}}Extrapolation error\\(not a gap)\end{tabular}}
 & MLP, 5 seeds & $+10.29\%\pm7.02$pp \\
 & LSTM, 5 seeds & $+16.18\%\pm9.39$pp \\
\midrule
Reference & Realized return & $+9.4\%$ \\
\bottomrule
\end{tabular}
\end{table}

\subsection{Reading the SCM weights: China exposure}\label{sec:exposure}

Why did the equal-$V$ SCM load 36\% on XLK? Not tariff exposure: XLK's
China-exposure proxy ($\sim$20\%) exceeds the index's ($\sim$7.5\%), and
the weighted synthetic exposure ($\sim$11.3\%) is \emph{higher} than the
actual index's. The weights reflect pre-period path coincidence (the AI
rally), which is exactly the homogeneity mechanism of
Section~\ref{sec:exxlk}. The exposure lens clarifies the weights'
economic content and nothing about any gap.

% ===========================================================================
\section{A Pre-Reporting Diagnostic Protocol, with a Positive Control}
\label{sec:protocol}

The tariff analysis is a failure case: every diagnostic fired. For the
diagnostics to be a protocol rather than a post-mortem, a reader must also
see what \emph{passing} looks like. This section states the protocol and
runs it end-to-end on the canonical clean-control application---California
Proposition~99 in \citet{abadie2010}---as a positive control.

\subsection{The checklist}\label{sec:checklist}

Before any synthetic-control gap is reported as an effect for a
market-wide or cross-sectional shock:

\begin{enumerate}
  \item \textbf{Donor--treatment homogeneity.} Compute the pre-period
  correlation between the treated unit's path and each donor's path
  (returns or first differences, not levels). If donors are slices of the
  treated unit or correlations are extreme, the synthetic control will
  reconstruct the treated series and the gap is pinned near zero by
  construction. \emph{Tariff case:} max correlation 0.91 (XLK), donors
  are index components $\Rightarrow$ fail
  (Sections~\ref{sec:vsens}, \ref{sec:exxlk}).
  \item \textbf{Placebo power.} Run permutation placebos and ask whether
  the realized gap is \emph{extreme} among them. If the realized gap is
  the smallest in absolute value, the test has no power against any
  alternative, and neither significance nor insignificance is
  interpretable. \emph{Tariff case:} $p=100\%$, smallest of 12 units
  $\Rightarrow$ fail (Section~\ref{sec:placebo}).
  \item \textbf{Weight concentration and pre-trends.} Inspect the weight
  vector (nonzero count, top share) and the pre-trend statistic.
  Single-donor domination and pre-trend failure disqualify the
  construction. \emph{Tariff case:} overseas pool top weight 0.76
  (Canada), pre-trend $p=0.000$ $\Rightarrow$ fail
  (Section~\ref{sec:overseas}); note that a \emph{high} pre-trend
  $p$-value is equally uninformative under step-2 failure.
  \item \textbf{Binding-tightness sweep.} Sweep a regularization dial
  (augmented-SCM ridge penalty) from fully donor-bound to unbound. If the
  gap orders monotonically along the dial---from the convex bound to the
  raw return---it tracks the control construction, not the treatment.
  \emph{Tariff case:} $+0.68\%\to+9.40\%$ monotone in $\alpha$
  $\Rightarrow$ fail (Section~\ref{sec:ascm}).
\end{enumerate}

A gap may be reported as an effect only when step 2 shows power and step 4
is flat. The tariff events fail all four steps in known directions, which
is why Section~\ref{sec:results} reports a diagnosis rather than
estimates.

\subsection{Positive control: California Proposition 99}\label{sec:poscontrol}

We apply the same four diagnostics to the classic clean-control case:
California's 1988 tobacco program, with 38 untreated donor states and
annual cigarette sales per capita, 1970--2000
(\texttt{positive\_control.py}; the ADH predictor set via \texttt{pysyncon};
data mirrored alongside the script). This is an environment the protocol
is supposed to pass: donors are other states, not slices of California,
and a large published effect exists against which the diagnostics can be
calibrated. Table~\ref{tab:poscontrol} reports all four steps side by
side with the tariff case.

The convex SCM recovers the published effect: California's 2000 gap is
$-25.5$ packs per capita, within the range reported by \citet{abadie2010}
(roughly $-18$ to $-25$ across specifications). Around that estimate the
diagnostics behave as a protocol should:

\begin{itemize}
  \item \textbf{Homogeneity.} State-level co-movement is not trivial
  (max first-difference correlation 0.84, mean 0.58), so high correlation
  \emph{alone} is not disqualifying. What the tariff case has and this
  case lacks is \emph{containment}: no donor state is a component of
  California, whereas XLK (correlation 0.91) is a component of the
  treated index, and the synthetic control reconstructs it.
  \item \textbf{Placebo power.} On the naive terminal-gap statistic the
  test is still weak (California ranks 7th of 39, $p=0.179$)---a useful
  reminder that raw terminal gaps are noisy even in the canonical case.
  On the ADH-style post/pre RMSPE ratio, California is the single most
  extreme unit of 39 (ratio 9.2, permutation $p=0.026$). The contrast
  with the tariff case is photographic: there the realized gap is the
  \emph{smallest} of all units ($p=1.00$).
  \item \textbf{Concentration.} Weights spread over 9 of 38 donors
  (top-1 0.31, top-3 0.66), versus 0.76 on a single index in the tariff
  overseas pool.
  \item \textbf{Binding tightness.} The effect exists under the tightest
  binding (convex: $-25.5$); the ridge dial varies its magnitude within
  the same regime ($-19.5$ to $-33.8$ over $\alpha\in[0.01,100]$)
  without approaching the raw change ($-48.5$) or collapsing to zero.
  Here the dial adjusts the estimate; in the tariff case the dial
  \emph{was} the estimate.
\end{itemize}

\begin{table}[htbp]
\centering
\caption{The four diagnostics: tariff events (failure case) versus
California Proposition 99 (positive control). Positive-control numbers
from \texttt{positive\_control.py}.}
\label{tab:poscontrol}
\small
\setlength{\tabcolsep}{4pt}
\begin{tabular}{@{}p{2.9cm}p{5.4cm}p{5.7cm}@{}}
\toprule
Diagnostic & 2018/2025 tariff events & California Prop.\ 99 \\
\midrule
Donor homogeneity & max corr 0.91; donors are components of the treated
index & max corr 0.84, no containment (donors are other states; mean
0.58) \\
Placebo power & realized gap smallest of 12 units, $p=1.00$ & RMSPE-ratio
most extreme of 39 units, $p=0.026$ (terminal-gap $p=0.179$) \\
Weight concentration & overseas top-1 0.76 (Canada) & 9 of 38 nonzero,
top-1 0.31, top-3 0.66 \\
Binding-tightness sweep & monotone: convex $+0.68\%$ $\to$ raw
$+9.40\%$ & effect present at every $\alpha$ ($-19.5$ to $-33.8$);
convex $-25.5$; raw change $-48.5$ \\
\bottomrule
\end{tabular}
\end{table}

% ===========================================================================
\section{Discussion}\label{sec:discussion}

\subsection{One contaminated world, many fractured estimands}
\label{sec:estimands}

The classical SCM application (California tobacco vs.\ other states)
assumes untreated donors. A global tariff shock violates this for every
financial asset simultaneously, in the same direction: all controls fall
with the treated unit, pulling every synthetic path down and compressing
every treated-minus-synthetic statistic toward zero relative to whatever
the unobservable counterfactual would have been. The dispersion across
specifications in Table~\ref{tab:map} measures how differently the various
constructions propagate that contamination---it is not uncertainty about
one quantity. A synthetic-control statistic in this setting is best
understood as a well-defined linear functional of (treated, donors)
paths whose causal content is nil.

\subsection{Why market adjustment cannot isolate sector effects here}
\label{sec:es-discussion}

In firm-level event studies the market model removes common noise because
the event is firm-specific. In a market-wide treatment the market factor
is the treatment, and each sector's own return sits inside it
(XLK is the largest weight in the S\&P~500). Residuals therefore mix
genuine relative strength with mechanical own-weight absorption.
Cross-sectional statements of the form ``sector X was abnormally
resilient'' are unavailable in this design; what is available is the
descriptive decomposition of rankings into beta and residual components.

\subsection{Two rounds, three interpretations}\label{sec:learning}

The volatility contrast (immediate spike, one-day half-life in 2025; no
immediate reaction in 2018) admits at least three readings: the market
learned the escalation template between rounds; the 2025 treatment was
larger and more surprising than any single 2018 node; the background
(rates, inflation, AI cycle) differed. Treatment intensity differs by an
order of magnitude across the focal dates, so nothing identifies the
learning channel. We present the contrast as fact and the interpretations
as open.

\subsection{Limitations}\label{sec:limitations}

\begin{enumerate}
  \item \textbf{No clean control}: every candidate donor is treated;
  this is the paper's subject and its boundary.
  \item \textbf{Survivorship}: \texttt{yfinance} covers listed tickers;
  mild for sector ETFs, nonzero.
  \item \textbf{2018 simplification}: one focal node represents a
  seven-node escalation.
  \item \textbf{VIX mechanical identity}: any return-on-$\Delta$VIX
  exercise is contemporaneously mechanical; covariation shares are
  descriptive only, and $\hat\beta$ is extrapolated across a regime
  shift.
  \item \textbf{$V$--$w$ coupling}: the data-driven predictor weights are
  chosen on pre-period outcomes and then interact with the weight
  estimation; optimizing pre-fit can worsen post-period validity. This
  known SCM critique is visible here as the 11$\times$ spread of
  Section~\ref{sec:vsens}.
  \item \textbf{Estimand discipline}: no statistic in
  Table~\ref{tab:map} is an effect; we have flagged each block's actual
  content.
  \item \textbf{China exposure is a proxy} (PIIE/BEA synthesis), used
  only to interpret weights.
  \item \textbf{Learning not identified}: treatment magnitude, surprise,
  and background all differ across rounds.
\end{enumerate}

% ===========================================================================
\section{Conclusion}\label{sec:conclusion}

This paper asked what counterfactual methods deliver when a policy shock
treats every candidate control, and answered with a diagnosis. Four
findings define the limits. Placebo tests have no power when donors are
homogeneous slices of the treated unit; convex SCM concentrates overseas
weights on a single best tracker and fails pre-trends; counterfactual
statistics order monotonically with a binding-tightness dial (the
augmented-SCM ridge sweep) that has nothing to do with the treatment; and
the apparent spread of ``estimates'' across methods is a collection of
different estimands, none causal. A mediation-style decomposition of the
crash through VIX was shown to be mechanically flawed---VIX is the
implied volatility of the treated index's own options---and was retained
only as descriptive covariation accounting.

The firm empirical content that survives is donor-free and descriptive:
around the 2025 Liberation Day announcement, VIX rose about 140\% from
its pre-announcement close to its April 8 peak and decayed with a
one-day half-life after the April~9 pause,
while the comparable 2018 node produced no immediate VIX reaction; and
the sector cross-section's raw decline rankings are dominated by market
beta. The 2026 legal collapse of the tariff authorities (IEEPA struck
down; Section~122 expired under injunction; Section~301 pending)
underlines the environment these markets were pricing: policy as a
moving target.

For practitioners, Section~\ref{sec:protocol} distills the diagnosis into
a four-step pre-reporting protocol---homogeneity, placebo power, weight
concentration, binding-tightness sweep---calibrated on a clean-control
positive control. A gap may be reported as an effect only when placebos
have power and the statistic is invariant to the dial. In the two tariff
rounds studied here---and, we suspect, in many applications to
synchronized market shocks---those conditions fail, and the responsible
output is the diagnosis itself.

% ===========================================================================
\bibliographystyle{plainnat}
\bibliography{references}

\appendix
\section{Reproducibility}\label{app:repro}

All results are generated by fifteen Jupyter notebooks, each produced by a
\texttt{generate\_0X\_*.py} script and sharing the library
\texttt{paper\_lib.py} (SCM and augmented SCM via \texttt{cvxpy}, BSTS via
\texttt{statsmodels}, event study, covariation accounting, placebo/LOO,
GARCH via \texttt{arch}). Data are pulled by
\texttt{download\_tariff\_data.py} (with
\texttt{download\_overseas\_standalone.py} as a rate-limit fallback). The
augmented-SCM sweep of Table~\ref{tab:ascm} is reproduced by
\texttt{python ascm\_sweep.py}; the positive control of
Section~\ref{sec:poscontrol} by \texttt{python positive\_control.py}
(California panel mirrored as \texttt{smoking\_california.csv}; the 39
placebo fits run once and are cached, so re-runs are instant); the
aligned-horizon comparison of Section~\ref{sec:tworound} by
\texttt{python aligned\_windows.py}. An
auxiliary GDELT news panel (1{,}416
items, 2025-04-01--15; notebook \texttt{nlp\_01}) was explored for a
sentiment dimension; its finance-specific tone proved weak and it is not
part of the analysis. Table~\ref{tab:notebooks} maps sections to
notebooks.

\begin{table}[htbp]
\centering
\caption{Notebook map}
\label{tab:notebooks}
\small
\begin{tabular}{@{}ll@{}}
\toprule
Notebook & Section \\
\midrule
01--02 data exploration / event EDA & \ref{sec:narrative} \\
03 sector SCM + $V$ sensitivity & \ref{sec:vsens} \\
04 CausalImpact + in-time placebo & \ref{sec:overseas} \\
05 event study residuals & \ref{sec:cross-section} \\
06 placebo / LOO & \ref{sec:placebo} \\
07 overseas SCM (discarded) & \ref{sec:overseas} \\
08 two-round comparison & \ref{sec:tworound} \\
09 retrospective & \ref{sec:results} (overview) \\
10 VIX orthogonalization & \ref{sec:covar} \\
11 GARCH regimes & \ref{sec:garch} \\
12 MLP/LSTM extrapolation errors & \ref{sec:dl} \\
13 covariation decomposition & \ref{sec:covar} \\
14 China exposure & \ref{sec:exposure} \\
15 exclude-XLK mechanism & \ref{sec:exxlk} \\
\bottomrule
\end{tabular}
\end{table}

\end{document}
